\chapter{Host Company}

\section{Profile and services}

CODELOGIC is a consulting firm specialized in software solutions [1]. It was
founded in 2023. Its services include custom software development, database
architecture, cloud solutions, security services, mobile development, and web
applications. The company also works on cloud and DevOps activities involving
continuous integration and delivery, scaling, and automation, as well as data
and API services involving secure integrations and data pipelines.

These activities cover several layers of software delivery. Custom software
development begins from a defined need and produces a solution for that need.
Database architecture organizes persistent information and the relationships
through which applications retrieve it. Cloud solutions provide computing and
storage resources. DevOps practices connect source control, automated checks,
deployment, and operation. Security services concern the protection of
software and information. Mobile and web applications provide interfaces for
users. Data pipelines and APIs connect sources and services while controlling
how information moves between them.

My internship project was a Python and Flask data-analysis prototype connected
to the company's data, API, and web-application activities. It read stored
financial-market files, applied a sequence of analytical operations, generated
review files, and presented the results through a browser dashboard. The
assignment therefore combined software development with statistical analysis
and data management. The original prototype's browser dashboard was later
replaced by the React Terminal in the final project.

The project required a clear separation between source data, derived results,
and presentation. Market and news records remained protected inputs. Python
modules produced the statistical indicators and review tables. Flask endpoints
then supplied the generated information to the dashboard. This organization
made each result traceable and kept the analytical rules separate from the
interface.

Company supervisor: Firas Jebari.

\clearpage

\section{Internship setting and professional experience}

My internship experience was organized around a complete technical assignment
rather than an isolated programming exercise. I worked with a local codebase,
stored financial records, generated analysis files, application endpoints, and
a browser interface. This required me to understand the purpose of each layer
and the way information moved from an input row to a result displayed for
review.

The work began with data preparation and validation. I checked the structure
of market and news records, preserved the source files, and organized derived
outputs so that each processing stage could be inspected. I then worked on the
statistical calculations used to identify unusual price returns and trading
volumes. These calculations required attention to chronological order,
rolling references, missing values, and long gaps between observations.

The next part of the assignment connected the statistical results with stored
news context and a narrower watchlist. I had to keep the signal, its context,
and the final selection rule separate. This distinction was important because
an unusual market observation is not automatically explained by a nearby
headline, and a headline does not prove the cause of a movement.

I also worked on presenting the results through Flask endpoints and dashboard
pages. The interface had to preserve the analytical details while remaining
clear enough for review. A user could inspect a flagged row, its market
history, its statistical scores, the available news context, and its watchlist
status. Automated tests and repeated output checks were used to verify that the
presentation remained consistent with the generated files.

This experience strengthened my autonomy, documentation practice, and
attention to traceability. I learned to record assumptions beside
calculations, follow a displayed value back to its source fields, and
investigate unexpected output before changing a rule. The project also
improved my ability to connect data preparation, statistical reasoning,
backend development, interface design, testing, and technical writing in one
coherent deliverable.
